\documentclass[12 pt]{article}

\usepackage{stoversymb}
\usepackage[margin=0.75in, top=0.875in, bottom=0.875in]{geometry}
\usepackage{amsmath,amsfonts,amssymb,url,multicol,graphicx,tikz,soul,changepage}
%===makes urls render well===
\usepackage{lmodern}
\usepackage[T1]{fontenc}
%============================
\usepackage{wasysym} % smileys
\usepackage[inline]{enumitem}
\usepackage[many]{tcolorbox}
\tcbuselibrary{breakable}

%\everymath{\displaystyle}

\makeatletter
\renewcommand*\env@matrix[1][\arraystretch]{%
	\edef\arraystretch{#1}%
	\hskip -\arraycolsep
	\let\@ifnextchar\new@ifnextchar
	\array{*\c@MaxMatrixCols c}}
\makeatother

%\setenumerate{itemsep=0.25in}
%\setlist[enumerate,1]{label=\arabic*., itemsep=0.5in}
%\setlist[enumerate,2]{label={(\alph*)}, itemsep=0.8in}
%\setlist[enumerate,2]{leftmargin=0.5in}
%\setlist[enumerate,3]{label={\roman*.}}
\setlist[itemize,1]{label=$\circ$,topsep=0mm,itemsep=0mm}

\newcommand{\truefalse}[1]{#1\hfill\rule[-1mm]{220pt}{0.75pt}}
\newcommand{\hint}[1]{\vspace{3mm}\textbf{Hint}: #1}
\newcommand{\note}[1]{\textbf{Note}: #1}
\newcommand{\infsum}[3]{\sum_{{#1}={#2}}^\infty {#3}}
\newcommand{\axes}[1]{\begin{center}\includegraphics[scale=#1]{3DAxes}\end{center}}
\newcommand{\comps}[1]{\langle #1_1,#1_2,#1_3\rangle}
\newcommand{\compslong}[3]{\langle #1, #2, #3\rangle}
\newcommand{\ijk}[2]{#1\vect{#2}}
%\newcommand{\pmat}[1]{\begin{pmatrix}[1.25]#1\end{pmatrix}}
%\newcommand{\bmat}[1]{\begin{bmatrix}[1.25]#1\end{bmatrix}}
\newcounter{Exx}
\newcommand{\exbox}[1]{\stepcounter{Exx}\begin{tcolorbox}[breakable,arc=2pt,boxrule=1pt,right=6mm,top=3mm,bottom=6mm]\noindent\textit{\large\bfseries Example \theExx:\\[3mm]}~#1\end{tcolorbox}}
\newcommand{\exit}[1]{\vspace{4.5mm}\noindent\textit{Example #1:}}
\newcommand{\exbf}[1]{\vspace{1.5mm}\noindent\textbf{Example #1:}}
\newcommand{\LHS}{\text{LHS}}
\newcommand{\RHS}{\text{RHS}}
\newcommand{\ans}{\noindent\underline{Ans:}~}

\newenvironment{mypmatrix}[1]{\renewcommand{\arraystretch}{#1}\begin{pmatrix}}{\end{pmatrix}}
\newenvironment{mybmatrix}[1]{\renewcommand{\arraystretch}{#1}\begin{bmatrix}}{\end{bmatrix}}
\newcommand{\justify}{\ul{Justify your claim.}}

\newcommand{\pmat}[1]{\begin{mypmatrix}{1.25}#1\end{mypmatrix}}
\newcommand{\bmat}[1]{\begin{mybmatrix}{1.25}#1\end{mybmatrix}}

\makeatletter
\renewcommand*\env@matrix[1][*\c@MaxMatrixCols c]{%
	\hskip -\arraycolsep
	\let\@ifnextchar\new@ifnextchar
	\array{#1}}
\makeatother

\newcommand{\pmatgrid}[2]{\renewcommand{\arraystretch}{1.25}\begin{pmatrix}[#1] #2\end{pmatrix}}
\newcommand{\bmatgrid}[2]{\renewcommand{\arraystretch}{1.25}\begin{bmatrix}[#1] #2\end{bmatrix}}
\newcommand{\warnbox}[1]{\begin{adjustwidth}{0.375in}{0.375in}\begin{tcolorbox}[arc=0pt, colback=red!5!white, colframe=red, boxrule=0.5pt, parbox=false]{\parskip=3mm \parindent=0mm \red{\textbf{\large\itshape Warning!}\par #1}}\end{tcolorbox}\end{adjustwidth}}
\newcommand{\greentheorembox}[2]{\begin{adjustwidth}{0.375in}{0.375in}\begin{tcolorbox}[arc=0pt, colback=darkgreen!10!white, colframe=darkgreen, boxrule=0.5pt, parbox=false]{\parskip=3mm \parindent=0mm \green{\textbf{\large\itshape #1:}\par #2}}\end{tcolorbox}\end{adjustwidth}}
\newcounter{FactBox}
\newcommand{\factbox}[1]{\stepcounter{FactBox}\begin{adjustwidth}{0.375in}{0.375in}\begin{tcolorbox}[arc=0pt, colback=darkgreen!10!white, colframe=darkgreen, boxrule=0.5pt, parbox=false]{\parskip=3mm \parindent=0mm \green{\textbf{\large\itshape Fact \theFactBox:}\par #1}}\end{tcolorbox}\end{adjustwidth}}

\newcommand{\arrow}[1]{\xrightarrow{\hspace*{#1}}}

\setlength\parskip{0.75mm}	

\newcounter{ThmBox}
\newcommand{\thmbox}[1]{\stepcounter{ThmBox}\vspace{6mm}\noindent\fbox{\parbox{\linewidth}{~~(\theThmBox)~\hfill #1 \hfill\hfill}}}
\newcommand{\why}[1]{\begin{adjustwidth}{0.5in}{0.5in}{\par\vspace{3mm}\noindent\green{\textbf{Why should this be true? }#1}}\qed\end{adjustwidth}\vspace{6mm}}

\graphicspath{ {./../img/} }
\DeclareGraphicsExtensions{.pdf}

\addtocounter{ThmBox}{9}

\begin{document}
%\begin{flushright}Name: \line(1,0){200}\end{flushright}
\begin{center}
\Large{\textbf{The Inverse Matrix Theorem II}}
\end{center}
%Recall that the \ul{inverse} of an $n\times n$ matrix $\sansA$ is an $n\times n$ matrix $\sansA^{-1}$ for which 
%\[
%	\sansA\sansA^{-1}=\sansI_n=\sansA^{-1}\sansA,
%\] 
%where $\sansI_n$ is the $n\times n$ identity matrix. Not all matrices have inverses, and those that do are called \ul{invertible} or \ul{nonsingular}.

As we saw before, a matrix being invertible/nonsingular tells you a \textit{tremendous amount} about the matrix + its corresponding linear system(s) / linear transformation(s). To highlight this, your textbook regularly adds to + revisits a monolithic colossus it calls \textbf{The Invertible Matrix Theorem}. 

This so-called ``theorem'' is \ul{really} just a collection of statements/observations which mean the same thing as (and hence are logically equivalent to) \textit{$\sansA$ has an inverse}. However, because many of the statements lumped into this ``theorem'' are important---and indeed, many are related to / duplicates of statements we've already visited before---I want to make sure you have them explicitly given and explained to you. Hence, this (and the previous) handout!

Here, we recall the nine (9) previously-stated conditions which are equivalent to the statement ``the $n\times n$ matrix $\sansA$ is invertible:''
\begin{enumerate}[label=(\arabic*),itemsep=0.75mm,topsep=1.5mm,leftmargin=16mm]
	\item $\sansA$ is invertible if and only if $\det(\sansA)\neq 0$.
	\item $\sansA$ is invertible if and only if the columns of $\sansA$ form a linearly independent set.
	\item $\sansA$ is invertible if and only if $\sansA\vect{x}=\vect{0}$ has only the trivial solution.
	\item $\sansA$ is invertible if and only if the linear transformation $T(\vect{x})=\sansA\vect{x}$ is one-to-one.
	\item $\sansA$ is invertible if and only if the linear transformation $T(\vect{x})=\sansA\vect{x}$ is onto.
	\item $\sansA$ is invertible if and only if the equation $\sansA\vect{x}=\vect{b}$ has $\geq 1$ solution for each $\vect{b}\in\Reals^n$.
	\item $\sansA$ is invertible if and only if the columns of $\sansA$ span $\Reals^n$.
	\item $\sansA$ is invertible if and only if $\sansA$ is row equivalent to the $n\times n$ identity matrix $\sansI_n$.
	\item $\sansA$ is invertible if and only if $\sansA^\sansT$ invertible.
\end{enumerate}
For more details, you should see the previous handout.

With this handout, we pick up from there. \red{\textbf{Throughout, assume that $\sansA=\pmatgrid{c|c|c}{\vect{v}_1 & \cdots & \vect{v}_n}$ is an $n\times n$ matrix, i.e. that $T(\vect{x})=\sansA\vect{x}$ is a linear transformation $\Reals^n\to\Reals^n$!}}


%==========================================================================
\thmbox{$\sansA$ is invertible if and only if the columns of $\sansA$ form a basis for $\Reals^n$.}
\why{From the previous handout, the $n\times n$ matrix $\sansA$ is invertible if and only if the columns of $\sansA$ are linearly independent (see (2) on the last handout) \textit{and} if and only if the columns of $\sansA$ span $\Reals^n$ (see (7) on the last handout).

To rephrase: $\sansA=\pmatgrid{c|c|c}{\vect{v}_1 & \cdots & \vect{v}_n}$ is invertible if and only if $\{\vect{v}_1,\ldots,\vect{v}_n\}$ is a linearly independent set which spans $\Reals^n$. By definition, this means that $\{\vect{v}_1,\ldots,\vect{v}_n\}$ is a basis for $\Reals^n$.
}

%==========================================================================
Condition (10) involves the columns of $\sansA$; unsurprisingly, we can relate this to the column \textit{space} $\col(\sansA)$ of $\sansA$.

\thmbox{$\sansA$ is invertible if and only if $\underbrace{\col(\sansA)}_{\row(\sansA^\sansT)}=\Reals^n$.}
\why{The columns $\{\vect{v}_1,\ldots,\vect{v}_n\}$ of $\sansA$ are a basis for $\Reals^n$ if and only if $\col(\sansA)\deff\vsspan\{\vect{v}_1,\ldots,\vect{v}_n\}$ is an $n$-dimensional subspace of $\Reals^n$. 
	
By properties of (vector) subspaces, the only such subspace \textit{in} $\Reals^n$ is $\Reals^n$ itself.}

%==========================================================================
\thmbox{$\sansA$ is invertible if and only if $\dim(\underbrace{\col(\sansA)}_{\row(\sansA^\sansT)})=n$.}
\why{By (11), $\sansA$ is invertible if and only if $\col(\sansA)=\Reals^n$, which is true if and only if $\dim(\col(\sansA))=\dim(\Reals^n)=n$.}

%==========================================================================
\thmbox{$\sansA$ is invertible if and only if $\underbrace{\rank(\sansA)}_{\rank(\sansA^\sansT)}=n$.}
\why{By (12), $\sansA$ is invertible if and only if $\dim(\col(\sansA))=n$. By definition, $\rank(\sansA)=\dim(\col(\sansA))$, so the result follows.}
	
%==========================================================================
Conditions (11)--(13) involve $\col(\sansA)$, and by the rank-nullity theorem, this is immediately tethered to the null space $\nul(\sansA)$.

\thmbox{$\sansA$ is invertible if and only if $\nullity(\sansA)=0$}
\why{By (13), $\sansA$ is invertible if and only if $\rank(\sansA)=n$. By the rank-nullity theorem,
\[
	\rank(\sansA)+\nullity(\sansA)=n,
\]
so then $\sansA$ is invertible if and only if $\nullity(\sansA)=n-\rank(\sansA)=n-n=0$.}

%==========================================================================
\thmbox{$\sansA$ is invertible if and only if $\dim(\nul(\sansA))=0$}
\why{By (14), $\sansA$ is invertible if and only if $\dim(\nul(\sansA))=0$. By definition, $\nullity(\sansA)=\dim(\nul(\sansA))$, so the result follows.}

%==========================================================================
\thmbox{$\sansA$ is invertible if and only if $\nul(\sansA)=\{\vect{0}\}$.}
\why{By (15), $\sansA$ is invertible if and only if $\dim(\nul(\sansA))=0$. The only $0$-dimensional subspace of $\Reals^n$ (or of any vector space) is the trivial subspace $\{\vect{0}\}$, and hence the result follows.}

%==========================================================================
\vspace{3mm}\noindent Even when the two above lists are \textbf{combined}, the result is still just a small fraction of the number of equivalent ways one can say ``$\sansA$ is invertible.'' In a perfect world, we'll revisit this handout with additional updates at least one more time throughout the remainder of the semester, but this is less than certain.

In the meantime, take some time to read through the above and digest everything thoroughly. To help, use the above to work through the following true/false practice problems!

\vspace{3mm}

\exbox{Mark each of the following questions ``true'' or ``false.'' Throughout, let $\vect{v}_1,\ldots,\vect{v}_p$ be vectors in a nonzero subspace $H$ of $\Reals^n$ and let $S=\{\vect{v}_1,\ldots,\vect{v}_p\}$. \justify
\begin{enumerate}[label=(\alph*)]
	\item The set of all linear combinations of $\vect{v}_1,\ldots,\vect{v}_p$ is a subspace of $\Reals^n$.
	\item If $\{\vect{v}_1,\ldots,\vect{v}_{p-1}\}$ spans $H$, then $S$ spans $H$.
	\item If $\{\vect{v}_1,\ldots,\vect{v}_{p-1}\}$ is linearly independent, then so is $S$.
	\item If $S$ is linearly independent, then $S$ is a basis for $H$.
	\item If $\vsspan\{S\}=H$, then some subset of $S$ is a basis for $H$.
	\item If $\dim H=p$ and $\vsspan\{S\}=H$, then $S$ cannot be linearly dependent.
	\item A plane in $\Reals^3$ is a two-dimensional subspace.
	\item Row operations on a matrix $\sansA$ can change the linear dependence relations among the rows of $\sansA$.
	\item Row operations on a matrix can change the null space.
	\item The rank of a matrix equals the number of nonzero rows.
	\item If an $m\times n$ matrix $\sansA$ is row equivalent to an echelon matrix $\sansU$ and if $\sansU$ has $k$ nonzero rows, then the dimension of the solution space of $\sansA\vect{x}=\vect{0}$ is $m-k$.
	\item If $\sansB$ is obtained from $\sansA$ by elementary row operations, then $\rank(\sansB)=\rank(\sansA)$.
	\item The nonzero rows of a matrix $\sansA$ form a basis for $\row(\sansA)$.
	\item If matrices $\sansA$ and $\sansB$ have the same $\RREF$, then $\row(\sansA)=\row(\sansB)$.
	\item If $H$ is a subspace of $\Reals^3$, then there is a $3\times 3$ matrix $\sansA$ such that $H=\col(\sansA)$.
	\item If $\sansA$ is $m\times n$ and $\rank(\sansA)=m$, then the linear transformation $\vect{x}\mapsto\sansA\vect{x}$ is one-to-one.
	\item If $\sansA$ is $m\times n$ and the linear transformation $\vect{x}\mapsto\sansA\vect{x}$ is onto, then $\rank(\sansA)=m$.
\end{enumerate}
}

\vspace{3mm}

\exbox{What would you have to know about the solution set of a homogoenous system of 31 linear equations in 33 variables in order to know that the associated nonhomogeneous equation has a solution?}

\vspace{3mm}

\exbox{Let $\rmTT:\Reals^n\to\Reals^m$ be a linear transformation.
	\begin{enumerate}[label=(\alph*)]
		\item What is the dimension of $\range(\rmTT)$ if $\rmTT$ is one-to-one? \justify
		\item What is the dimension of $\ker(\rmTT)$ if $\rmTT$ is onto? \justify
	\end{enumerate}
}

\vspace{3mm}

\exbox{
	\red{This is a challenge problem!} Let $\sansA$ be an $m\times n$ matrix.
	
	\begin{enumerate}[label=(\alph*),listparindent=0.25in]
		\item Show that if $\sansB$ is $n\times p$, then $\rank(\sansA\sansB)\leq\rank(\sansA)$ by showing that every vector in the column space of $\sansA\sansB$ is in the column space of $\sansA$.
		\item Show that if $\sansB$ is $n\times p$, then $\rank(\sansA\sansB)\leq\rank(\sansB)$ by using problem 4(a) and studying $(\sansA\sansB)^\sansT)$.
		\item Show that if $\sansP$ is an invertible $m\times m$ matrix, then $\rank(\sansP\sansA)=\rank(\sansA)$ by applying problems 4(a) and 4(b) to each of $\sansP\sansA$ and $\sansP^{-1}(\sansP\sansA)$.
		\item Show that if $\sansQ$ is invertible, then $\rank(\sansA\sansQ)=\rank(\sansA)$ by applying problem 4(c) to $\rank(\sansA\sansQ)^\sansT$.
		\item Suppose that $\sansB$ is $n\times p$ such that $\sansA\sansB=0$. Show that $\rank(\sansA)+\rank(\sansB)\leq n$ by showing that one of $\nul(\sansA)$, $\col(\sansA)$, $\nul(\sansB)$, or $\col(\sansB)$ is contained in one of the other three.
		\item Suppose that $\rank(\sansA)=r$. The \textbf{rank factorization} of $\sansA$ is an equation of the form $\sansA=\sansC\sansR$ where $\sansC$ is an $m\times r$ matrix of rank $r$ and $\sansR$ is an $r\times n$ matrix of rank $r$. Such a factorization always exists.
		
		Given an $m\times n$ matrix $\sansB$, use rank factorizations of $\sansA$ and $\sansB$ to show that%\vspace{-1.5mm}
		\[
			\rank(\sansA+\sansB)\leq\rank(\sansA)+\rank(\sansB) 
		\]
		\vspace{-2mm}by writing the sum $\sansA+\sansB$ as the product of two augmented matrices.
		\item If $\sansA$ has rank $r$, explain why 
		\begin{enumerate}[label=(\roman*),topsep=0mm]
			\item $\sansA$ must contain an $m\times r$ submatrix $\sansA_1$ of rank $r$; and
			\item $\sansA_1$ must have an invertible $r\times r$ submatrix $\sansA_2$.
		\end{enumerate}
	\end{enumerate}
}
\end{document}